% !TEX root = ../../hw05.tex
\begin{proof}[Solution]
  Let $\varepsilon>0$. Since $f$ is continuous at $x_0$, there is a
  $\delta>0$ such that
  \[
    |x-x_0|<\delta
      \quad\Longrightarrow\quad |f(x)-f(x_0)|<\varepsilon.
  \]
  We can use the same $\delta$ for $g(x)=|f(x)|$. The reason is the
  reverse triangle inequality,
  \[
    \bigl||u|-|v|\bigr|\le |u-v|,\qquad u,v\in\R.
  \]
  Taking absolute values cannot increase the distance between two
  outputs, as illustrated in \autoref{fig:hw05:absolute-value-fold}.

  \begin{marginfigure}
    \centering
    % !TEX root = ../../problems/hw05.tex
% The two number lines use the same scale; u < 0 < v < |u|.
% Shared physical dimensions for interval and number-line diagrams.
\tikzset{
  interval diagram/.style={
    x=1cm,y=6mm,font=\footnotesize,
    color=black,line width=0.5pt,line cap=round,line join=round,
    >={Stealth[length=4pt,width=3pt]}
  },
  interval endpoint/.style={
    circle,draw=black,line width=0.5pt,
    minimum size=4pt,inner sep=0pt,outer sep=0pt
  },
  interval open/.style={interval endpoint,fill=white},
  interval closed/.style={interval endpoint,fill=black},
  interval label/.style={inner sep=1pt}
}

\begin{tikzpicture}[interval diagram,x=8mm,y=1mm]
  \draw[<->,gray] (-2.4,0) -- (2.4,0);
  \foreach \position/\name in {-2/u,0/0,1/v}{
    \node[interval closed] at (\position,0) {};
    \node[below=5pt] at (\position,0) {$\name$};
  }
  \draw[<->] (-2,8) -- (1,8)
    node[midway,above=4pt] {$|u-v|$};

  \draw[->,gray] (-2,-7)
    .. controls (-2,-17) and (2,-16) .. (2,-25);
  \node[anchor=west] at (-2.35,-21) {fold at $0$};

  \draw[->,gray] (-0.3,-30) -- (2.4,-30);
  \foreach \position/\name in {0/0,1/{|v|},2/{|u|}}{
    \node[interval closed] at (\position,-30) {};
    \node[below=5pt] at (\position,-30) {$\name$};
  }
  \draw[<->] (1,-41) -- (2,-41)
    node[midway,below=4pt] {$\bigl||u|-|v|\bigr|$};
\end{tikzpicture}

    \caption[Absolute value as a fold of the real line.]{%
      Absolute value folds the negative half-line onto the
      positive half-line. In the illustrated case $u<0<v<|u|$, the
      original distance is $|u|+v$, while the folded distance is
      $|u|-v$. Folding brings the two points closer together.}
    \label{fig:hw05:absolute-value-fold}
  \end{marginfigure}

  Applying the inequality with $u=f(x)$ and $v=f(x_0)$ gives
  \[
    |g(x)-g(x_0)|
      =\bigl||f(x)|-|f(x_0)|\bigr|
      \le |f(x)-f(x_0)|<\varepsilon
  \]
  whenever $|x-x_0|<\delta$. Hence $g$ is continuous at $x_0$.
\end{proof}
